% Generated by make_register.py from hypotheses.md; do not edit by hand.
\begin{table*}[t]
\caption{Research log, hypotheses 1 to 12 of 21: every hypothesis, one row each, with the verdict.}
\label{tab:register-1}
\centering\scriptsize
\begin{tabular}{@{}p{0.9cm}p{0.9cm}p{5.0cm}p{7.6cm}p{1.9cm}@{}}
\toprule
ID & Sec. & Hypothesis & Result & Verdict \\
\midrule
O1 & V-A & The netlists of Odin III through VTR's ABC stage (leg O') route in VPR and are comparable in quality to the Parmys flow (leg P): geomean of O'/P at most 1.20 and no design over 1.50 on CLBs, channel width, delay and wirelength; num\_mult and num\_memories equal or the difference explained. (As first written O/P with the linked ABC; only leg O' was run, so the verdict is on O'/P.) & all 15 route; every gated geomean inside 1.20 (LUT 1.041, CLB 1.012, memories 1.000, mults 1.002, W 0.985, CPD 0.986, WL 1.034); worst per-design gated ratios 1.288 (stereovision3 WL), 1.233 (ch\_intrinsics WL); blob\_merge LUT 1.777 with CLB 1.038; num\_mult differs on stereovision1 (44 against 43, 1.023) and stereovision2 (194 against 192, 1.010), not explained in the table; hardblocks\_lost WARNING (adders with no reader after VTR's ABC): mkPktMerge 19, boundtop 90, mkSMAdapter4B 35, stereovision1 120, mkDelayWorker32B 706 (M2-post-fix qor.md) & \textsc{possibly good} \\
O2 & V-A & The Parmys leg reproduces VTR's own golden results, so it is a valid baseline. & routed 21 of 21; geomeans within 0.5\% of golden; BLIF byte-identical on 21 of 21; M1: the P' leg routes all 19 designs (Part II measured table) & \textsc{useful} \\
O3 & V-A & Odin III reaches per-design parity with Parmys. & geomean ratios 0.985 to 1.041 but per-design outliers (blob\_merge LUT 1.777, stereovision3 WL 1.288) & \textsc{inconclusive} \\
O4 & V-B & \code{lower} is a refinement of the word-level design on the micros. & 399 ok, 140 expected, 0 unexpected; sweep: 518 compared, 0 differ, 8 unproven at the SAT size limit & \textsc{useful} \\
O5 & V-F & Odin III's netlists are equivalent to Parmys's on the 14 small VTR designs. & Phase 2 exit: 9 of 14 proved, 5 untested at the SAT size limit; after \code{lower}: 8 ok, 1 empty, 5 unproven, 0 differ & \textsc{inconclusive} \\
O6 & V-C & Odin III's netlists behave like the RTL under 50,000 random vectors. & Odin III passes 12 of 15; 3 untested (raygentop, boundtop, stereovision0), not bugs; the same check fails the Parmys netlist on 8 of 15 & \textsc{useful} \\
O7 & V-D & Rerunning the exit flow gives byte-identical netlists. & 15 of 15 byte-identical & \textsc{useful} \\
O8 & V-E & The gate flow runs the VerilogEval references. & 144 elaborate; 12 refused with a pinned reason (2 explicit casts, 8 inferred latches, 2 compound bit-select assignments) & \textsc{useful} \\
O9 & V-G & Checking against oracles finds defects in the Parmys flow. & D2b, D4, D17, D19, \code{BUF\_NODE}, \code{proc\_mux} casez loss (6 listed), each with its evidence & \textsc{useful} \\
O10 & V-B & \code{lower} agrees with Yosys's \code{techmap} on random modules. & 200 seeds: 200 ok; 5,000 seeds: 4,830 ok, 142 structural, 28 empty, 0 bugs & \textsc{useful} \\
O11 & V-K & The hard-block rules reproduce Parmys's. & every count difference explained (fsm recoding, flip-flop-bit merge, wreduce order, opt\_share); none from a mapping rule & \textsc{possibly good} \\
O12 & V-K & The BLIF and in-memory ABC exchanges give equivalent results. & every finished design equivalent; LUT ratios 0.95 to 1.02, sha 1.077 & \textsc{possibly good} \\
\bottomrule
\end{tabular}
\end{table*}

\begin{table*}[t]
\caption{Research log, hypotheses 13 to 21 of 21: every hypothesis, one row each, with the verdict.}
\label{tab:register-2}
\centering\scriptsize
\begin{tabular}{@{}p{0.9cm}p{0.9cm}p{5.0cm}p{7.6cm}p{1.9cm}@{}}
\toprule
ID & Sec. & Hypothesis & Result & Verdict \\
\midrule
O14 & V-H & Provenance records the construct chain and the joins correctly, and leaves every writer's output unchanged. & 541 ctest cases, 0 fail (2 skip); lint all steps pass; writer outputs byte-identical across --prov-ast modes on 539 micros (8,927 files per side) & \textsc{useful} \\
O15 & V-H & The name map plus netlist can be read back to the same netlist (the read back, split from O14). & BLIF with map 68 of 68; JSON 4 of 68; Verilog 23 of 68 (24 ambiguous top) & \textsc{inconclusive} \\
O16 & V-I & \code{--prov-ast refs} costs at most 5\% read time and 10 MB RSS on five large designs (LU32PEEng, LU64PEEng, arm\_core, mcml, mkDelayWorker32B). & time ratios 1.037 (LU32PEEng), 0.998 (LU64PEEng), 1.067 (arm\_core), 1.029 (mcml), 1.054 (mkDelayWorker32B); RSS delta under 1 MB; margin: extra read time 1.8 ms (mkDelayWorker32B, 33.7 to 35.5 ms) and 2.7 ms (arm\_core, 39.1 to 41.8 ms), 0.4 and 1.7 percentage points over the bound, against off-mode run-to-run sd 1.2 and 1.5 ms (docs/V1.2.0-task4-gate.txt) & \textsc{dismissed} \\
O17 & V-J & The attribution joins VPR's timing report to HDL without changing VPR's result. & identical with and without the extra report; no unmatched atoms on 6 designs; (companion work, the nornsIII repository, docs/paper/r5-results.tex, github.com/drpaj12/nornsIII, Norns III round 5: the attribution agreed with VPR's CPD within 0.5\% in 999 of 999 runs) & \textsc{possibly good} \\
U1 & VII & The cost of provenance names in VPR is small: the names added to the netlist do not change VPR's result or run time. & not run; needs: Leg Oprov of the QoR harness (Part II, measured results): the same 15 designs through VPR with provenance names against without, comparing the four result metrics and run time. & \textsc{untested} \\
U2 & VII & Odin III also stays within the gate on the five largest VTR designs (\code{mcml}, \code{LU32PEEng}, \code{bgm}, \code{LU8PEEng}, \code{LU64PEEng}). & not run; needs: Run legs P and O' on these five on a machine with more memory, with the same gate. & \textsc{untested} \\
U3 & VII & Latches with reset and the \code{\$\_DFFSR\_*} cell types map correctly. & not run; needs: Add them to the micro set and run the refinement proof and the functional check. & \textsc{untested} \\
U4 & VII & The plain-LUT architecture on \code{mkDelayWorker32B} passes through ABC. & not run; needs: Rerun with a longer ABC limit than the 3,600 s that stopped it. & \textsc{untested} \\
U5 & VII & Near-critical paths (not only the critical path) can be attributed to HDL objects, so the Mimir dataset can say why a structure sits on a timing-critical or near-critical path. & not run; needs: Use the v1.1.1 timing attribution with a window, report per HDL construct the share of near-critical paths over the 15 designs, check that VPR's result is unchanged. & \textsc{untested} \\
\bottomrule
\end{tabular}
\end{table*}
